\section*{Bab \ref{ch:b1:nernst} --- Kesetimbangan Redoks dan Persamaan Nernst}

\begin{solution}{exo:b1:nernst:1}
N dalam \ce{NH3}: $-$III; dalam \ce{NO3-}: $+$V. S dalam \ce{SO4^2-}:
$+$VI. O dalam \ce{H2O2}: $-$I. Cr dalam \ce{Cr2O7^2-}: $+$VI. Cl dalam
\ce{ClO-}: $+$I. Fe dalam \ce{Fe3O4}: rata-rata $+8/3$ (satu besi(II) dan
dua besi(III)). C dalam \ce{CH3OH}: $x + 4(+1) + (-2) = 0$, $-$II.
\end{solution}

\begin{solution}{exo:b1:nernst:2}
\ce{Cr2O7^2- + 14H+ + 6e- -> 2Cr^3+ + 7H2O};
\ce{H2O2 + 2H+ + 2e- -> 2H2O};
\ce{O2 + 2H+ + 2e- -> H2O2};
\ce{SO4^2- + 4H+ + 2e- -> SO2 + 2H2O}.
\end{solution}

\begin{solution}{exo:b1:nernst:3}
$E = -0.76 + 0.030\log[\ce{Zn^2+}]$; $E = 1.36 + 0.030\log(p(\ce{Cl2})/[\ce{Cl-}]^2)$;
$E = 1.23 + 0.015\log(p(\ce{O2})[\ce{H+}]^4)$; $E = 0.27 - 0.059\log[\ce{Cl-}]$
(tekanan dalam bar). Dalam seng sulfat \qty{0.010}{mol/L}:
$E = -0.76 + 0.030 \times (-2) = -0.82$~V.
\end{solution}

\begin{solution}{exo:b1:nernst:4}
$E(\ce{Zn}) = -0.76 + 0.030\log 0.10 = -0.79$~V, $E(\ce{Cu}) = 0.34 +
0.030\log 0.010 = 0.28$~V. Tegangan $0.28 - (-0.79) = 1.07$~V. Seng, yang
potensialnya lebih rendah, adalah kutub negatif, tempat oksidasi
berlangsung: \omterm{def:b1:nernst:cell}{anode}.
\end{solution}

\begin{solution}{exo:b1:nernst:5}
\ce{MnO4- + 5Fe^2+ + 8H+ -> Mn^2+ + 5Fe^3+ + 4H2O}, $n = 5$,
$\log K = 5(1.51 - 0.77)/0.059 = 63$: \omterm{def:b1:extent-q-and-k:quantitative}{kuantitatif} dan cepat, sebuah
titrasi klasik (permanganat menjadi indikatornya sendiri).
\end{solution}

\begin{solution}{exo:b1:nernst:6}
Logam-logam yang $E^\circ$-nya di bawah 0 bereaksi: Mg, Zn, Fe, Ni, Pb
(timbal hanya perlahan: selisihnya kecil); Cu dan Ag tidak. \ce{Mg + 2H+ -> Mg^2+ + H2}, $\log K = 2 \times 2.36/0.059 = 80$;
\ce{Fe + 2H+ -> Fe^2+ + H2}, $\log K = 2 \times 0.41/0.059 = 13.9$.
\end{solution}

\begin{solution}{exo:b1:nernst:7}
$n = 2$, $\log K = 2(0.53 - 0.02)/0.059 = 17$: \omterm{def:b1:extent-q-and-k:quantitative}{kuantitatif}.
\end{solution}

\begin{solution}{exo:b1:nernst:8}
$3E^\circ(\ce{Fe^3+}/\ce{Fe}) = 2(-0.41) + 0.77$, $E^\circ = -0.02$~V.
\ce{2Fe^3+ + Fe -> 3Fe^2+}: \omterm{def:b1:nernst:couple}{oksidator} \ce{Fe^3+} (0.77) di atas \omterm{def:b1:nernst:couple}{reduktor}
\ce{Fe} ($-0.41$), $n = 2$, $\log K = 2(0.77 + 0.41)/0.059 = 40$. Serbuk
besi menjaga larutan besi(II) agar tidak berubah menjadi besi(III).
\end{solution}

\begin{solution}{exo:b1:nernst:9}
\ce{O2 + 4H+ + 4e- -> 2H2O}: $E = 1.23 + \frac{0.059}{4}\log [\ce{H+}]^4 =
1.23 - 0.059\,\mathrm{pH}$. \ce{2H+ + 2e- -> H2}: $E = \frac{0.059}{2}\log
[\ce{H+}]^2 = -0.059\,\mathrm{pH}$. Selisihnya $1.23$~V pada pH berapa pun:
\omterm{def:b1:nernst:cell}{tegangan sel} bahan bakar itu tidak bergantung pada pH.
\end{solution}

\begin{solution}{exo:b1:nernst:10}
$E^\circ(\ce{Cu+}/\ce{Cu}) = 0.52 > E^\circ(\ce{Cu^2+}/\ce{Cu+}) = 0.16$:
\omterm{def:b1:nernst:couple}{oksidator} \ce{Cu+} pasangan atas bereaksi dengan \omterm{def:b1:nernst:couple}{reduktor} \ce{Cu+}
pasangan bawah. $\log K = (0.52 - 0.16)/0.059 = 6.1$, $K =
[\ce{Cu^2+}]/[\ce{Cu+}]^2$, sehingga $[\ce{Cu+}] = \sqrt{0.10/10^{6.1}} =
\qty{2.8e-4}{mol/L}$.
\end{solution}

\begin{solution}{exo:b1:nernst:11}
$U = 0.059\log(0.10/\num{1.0e-3}) = 0.118$~V. Kutub positifnya adalah
perak dalam larutan yang pekat, tempat \ce{Ag+} tereduksi; dalam larutan
yang encer perak teroksidasi. Larutan yang pekat menjadi lebih encer dan
yang encer menjadi lebih pekat; tegangannya turun dan mencapai nol ketika
kedua konsentrasi sama.
\end{solution}

\begin{solution}{exo:b1:nernst:12}
\ce{H2O2} adalah \omterm{def:b1:nernst:couple}{oksidator} pasangan \ce{H2O2}/\ce{H2O} (1.76~V) dan
\omterm{def:b1:nernst:couple}{reduktor} pasangan \ce{O2}/\ce{H2O2} (0.69~V): zat ini bereaksi dengan
dirinya sendiri, \ce{2H2O2 -> 2H2O + O2},
$\log K = 2(1.76 - 0.69)/0.059 = 36$. Reaksi itu disukai tetapi sangat
lambat tanpa \omterm{def:b1:elementary-steps:catalyst}{katalis} (\cref{ch:b1:elementary-steps}); botol yang bersih
serta penyimpanan yang dingin dan gelap menjaganya tetap lambat.
\end{solution}

\begin{solution}{pb:b1:nernst:1}
\textbf{1.} Ag: 0, $+$I, $+$I; Cl: $-$I pada keduanya.
\textbf{2.} \ce{AgCl(s) + e- -> Ag(s) + Cl-}.
\textbf{3.} $E = E^\circ(\ce{AgCl}/\ce{Ag}) - 0.059\log[\ce{Cl-}]$.
\textbf{4.} Kedua padatan beraktivitas 1; hanya ion klorida yang berada
dalam larutan.
\textbf{5.} $E^\circ = 0.22$~V.
\textbf{6.} $0.22 + 0.118 = 0.34$~V.
\textbf{7.} \ce{Pt} $|$ \ce{H2} (\qty{1}{bar}) $|$ \ce{H+} (\omterm{def:b1:extent-q-and-k:activity}{aktivitas} 1)
$\|$ \ce{Cl-} (\qty{0.10}{mol/L}) $|$ \ce{AgCl} $|$ \ce{Ag}.
\textbf{8.} Elektrode perak: $0.22 + 0.059 = 0.28$~V $> 0$.
\textbf{9.} $0.28$~V.
\textbf{10.} \ce{2AgCl + H2 -> 2Ag + 2H+ + 2Cl-}.
\textbf{11.} $\log K = 2 \times 0.22/0.059 = 7.5$.
\textbf{12.} Dari platina (kutub negatif, tempat \ce{H2} teroksidasi) ke
perak.
\textbf{13.} $0.27 - 0.059\log 1.0 = 0.27$~V.
\textbf{14.} $E(\ce{Ag}) = 0.22 + 0.177 = 0.40$~V $> 0.27$: elektrode
perak positif, $U = 0.13$~V.
\textbf{15.} $U = 0.22 - 0.059\log[\ce{Cl-}] - 0.27 = -0.05 -
0.059\log[\ce{Cl-}]$ (volt).
\textbf{16.} $\log[\ce{Cl-}] = -(0.115 + 0.05)/0.059 = -2.80$,
$[\ce{Cl-}] = \qty{1.6e-3}{mol/L}$, yaitu \qty{57}{mg/L}.
\textbf{17.} $\Delta c/c = \ln 10 \times 0.001/0.059 = 0.039$: sekitar
\qty{4}{\percent} per milivolt.
\textbf{18.} Ketika klorida sampel menjadi sebanding dengan yang
dilepaskan oleh perak klorida itu sendiri, $\sqrt{K_s} = \qty{1.3e-5}{mol/L}$;
elektrode itu dipakai jauh di atasnya, mulai sekitar \qty{1e-4}{mol/L}.
\textbf{19.} $E = 0.80 + 0.059\log[\ce{Ag+}]$.
\textbf{20.} $[\ce{Ag+}] = K_s/[\ce{Cl-}]$.
\textbf{21.} $E = 0.80 + 0.059\log K_s - 0.059\log[\ce{Cl-}] = 0.80 -
0.059\,\mathrm{p}K_s - 0.059\log[\ce{Cl-}]$.
\textbf{22.} $E^\circ(\ce{AgCl}/\ce{Ag}) = E^\circ(\ce{Ag+}/\ce{Ag}) -
0.059\,\mathrm{p}K_s$ (\cref{prop:b1:nernst:second-kind}).
\textbf{23.} $0.80 - 0.059 \times 9.75 = 0.225$~V, sesuai dengan tabel
(0.22~V), yang dihitung secara terpisah dari energi Gibbs.
\textbf{24.} $E^\circ(\ce{AgCl}/\ce{Ag}) = $ \textbf{0.22~V}.
\end{solution}
